% 摘要正文：替换主文件中的占位摘要；沿用主文件的标题和摘要环境。
% 依据：当前《数学建模论文.pdf》正文及对应结果记录。

复杂交互场景中的情感分析需要综合文本、语音与视觉信息，而模态时间尺度不一致、局部信息缺失及预测依据不透明，增加了准确理解情感表达的难度。本文基于CMU-MOSEI相关数据，围绕多模态特征提取与时序对齐、缺失条件下的情感预测及预测结果解释，建立相应的数学模型与分析方法。

针对问题一，构建\textbf{基于词级时序对齐与时间交叠加权聚合的异构特征表征方法}。利用DistilBERT、openSMILE和MediaPipe分别提取上下文语义、声学低层描述量与面部表情系数。基于Wav2Vec2输出的字符概率进行连接时序分类（CTC）强制对齐，通过Viterbi解码估计候选词区间；结合音视频实际时间戳，按有效交叠时长归一化加权，将不同采样尺度的声学与视觉特征聚合至统一词位置。结合CTC置信度、诊断词错误率（WER）与人工复核进行质量筛查，并通过计算核验和固定窗口对照分析特征聚合与时序组织的合理性。最终输出覆盖100条样本、1932个词位置的三模态特征与原始素材映射，其中79.97\%的词通过对齐质量筛选，71.79\%同时满足筛选与三模态可观测条件。

针对问题二，构建\textbf{缺失感知门控融合的情感极性与强度联合预测模型}。采用附件2的对齐特征，通过不同模态组合、比例和位置的连续遮挡增强训练，结合残差时序卷积、掩码池化与观测质量统计，自适应分配模态融合权重，并利用分类与回归双输出分支完成预测。在727条测试样本的完整输入下，模型准确率为64.92\%，平均绝对误差（MAE）为0.6512；在三模态同步随机遮挡50\%有效位置的测试中，相较仅用完整输入训练的固定融合模型，宏平均F1由0.4328提高至0.5316，MAE由0.7936降至0.7304，体现了整体方案对局部连续缺失的适应能力。最终采用三模型集成，完成附件3全部30条样本的极性与强度预测。

针对问题三，建立\textbf{结合Shapley贡献与局部输入干预的分层解释方法}。在门控模型基础上引入模态组合增强，枚举八种模态组合精确计算贡献，以积分梯度筛选局部候选窗口，通过删除与保留实验检验其作用，并结合原文与语音对齐建立素材映射。在728条验证样本上，模型宏平均F1为0.6046，MAE为0.5651。按各样本绝对贡献最大的模态统计，删除候选窗口后，原预测类别概率平均下降12.84个百分点，高于等长随机窗口的4.64个百分点，表明候选片段更能反映模型所依赖的信息。对附件4全部20条样本输出预测、三模态贡献与关键证据，并标注时间定位状态。

本文将时序特征构建、局部缺失下的联合预测与关键证据回溯相结合，为复杂场景中的多模态情感分析及人工辅助审阅提供了可复核的方法与结果。

\par\noindent\textbf{关键词：}多模态情感分析；词级时序对齐；缺失感知门控融合；Shapley值；积分梯度
